\documentclass[t,10pt]{beamer}

% ==========================================
% EXACT CANVAS GEOMETRY
% ==========================================
\geometry{paperwidth=14.41cm, paperheight=9cm}

% ==========================================
% TYPOGRAPHY & MATH FONTS (Compile with XeLaTeX/LuaLaTeX)
% ==========================================
\usepackage{fontspec}
\usepackage[mathrm=sym]{unicode-math}

% Text Font: True Helvetica Clone
\setsansfont{Alegreya Sans}[
    Ligatures=TeX
]
\renewcommand{\familydefault}{\sfdefault}

% Math Font: New Computer Modern Sans
\setmathfont[
    Extension=.otf,
]{NewCMSansMath-Regular}

\setmonofont[
    Extension=.otf,
]{NewCMMono10-Regular}

% ==========================================
% COLORS & DESIGN
% ==========================================
\definecolor{ScreenshotBlue}{HTML}{2E6DB5}
\definecolor{TextBlack}{HTML}{000000}
\definecolor{SubtitleGray}{HTML}{333333}
\definecolor{BackgroundGray}{HTML}{F5F5F5}
\definecolor{BackgroundBlue}{HTML}{ECF2F9}

\setbeamertemplate{navigation symbols}{}
\setbeamercolor{background canvas}{bg=white}

% Margins matching the clean spacing of the handout
\setbeamersize{text margin left=0.3cm, text margin right=0.3cm}

% ==========================================
% EXACT MANUAL FONT SCALING (Bypassing Beamer defaults)
% ==========================================
% Regular body text size (Matches the handout text proportions)
\setbeamerfont{normal text}{size=\fontsize{11}{15}\selectfont}
\setbeamerfont{itemize/enumerate body}{size=\fontsize{11}{15}\selectfont}

% Slide Title size
\setbeamerfont{frametitle}{size=\fontsize{14}{18}\bfseries}

% Force our exact normal text size at document start
\AtBeginDocument{\usebeamerfont{normal text}}

% ==========================================
% CUSTOM TEMPLATES
% ==========================================
% Slide Titles
\setbeamercolor{frametitle}{fg=ScreenshotBlue}
\setbeamertemplate{frametitle}{
    \vspace*{0.2cm}
    {\usebeamerfont{frametitle}\insertframetitle}
    \vspace*{-0.2cm}
}

% Footline (Slide Numbering)
\setbeamertemplate{footline}{
    \hbox to \paperwidth{%
        \hfill
        \textcolor{SubtitleGray}{\fontsize{8}{10}\selectfont\insertframenumber}%
        \hspace*{0.3cm}%
    }
    \vskip0.3cm
}

\setbeamertemplate{frametopskip}{}

% Proportional Minimalist Square Bullets
\setbeamertemplate{itemize item}{\color{ScreenshotBlue}\raisebox{0.25ex}{\rule{0.6ex}{0.6ex}}}
\setbeamercolor{itemize item}{fg=ScreenshotBlue}
\setlength{\leftmargini}{1.2em}

\setbeamertemplate{enumerate item}{%
    \color{ScreenshotBlue}%
    \insertenumlabel.%
}
\setbeamercolor{enumerate item}{fg=ScreenshotBlue}
\AtBeginEnvironment{enumerate}{%
  \setlength{\leftmargini}{1.5em}% Adjust this value to get your perfect alignment
}

% Heading sizing for sub-sections inside frames
\newcommand{\slideheading}[1]{{\vspace*{0.3cm}\bfseries #1\par}}

\usepackage{tikz}
\usetikzlibrary{matrix, calc, positioning}
\usetikzlibrary{arrows.meta}
\tikzset{
  math node/.style={
    execute at begin node=$,
    execute at end node=$
  }
}

\usepackage[dvipsnames]{xcolor}

\usepackage{tcolorbox}
\tcbuselibrary{theorems}
\tcbuselibrary{skins}
\tcbuselibrary{hooks}

\let\definition\relax
\let\enddefinition\relax
\let\corollary\relax
\let\endcorollary\relax

\tcbset{
    thmstyle/.style={
	theorem style = plain, 
	enhanced jigsaw,
	fonttitle = \bfseries\color{ScreenshotBlue},
	%colback = BackgroundGray,         % Background color
	colback = BackgroundBlue,
	%colback = white,
	sharp corners,         % Make corners sharp
	%boxrule = 1pt,
	%colframe = ScreenshotBlue,
	frame hidden,
	before skip = 8pt,
	after skip = 7pt,
	%before upper app={\parindent15pt\ },
	%after=\noindent,
	top = 2mm,
	bottom = 2mm,
	left = 2mm, 
	right = 2mm
    }
}

\NewTcbTheorem{definition}{Definition}{thmstyle,colback = BackgroundGray}{def}
\NewTcbTheorem{proposition}{Proposition}{thmstyle}{prop}
\NewTcbTheorem{corollary}{Corollary}{thmstyle}{cor}

\let\proof\relax
\let\endproof\relax
\let\qedhere\relax

\newif\ifqedplaced

\NewDocumentEnvironment{proof}{O{Proof}}{%
  \global\qedplacedfalse
  \par
  \normalfont
  \setlength{\topsep}{0pt plus 0pt}%
  \setlength{\partopsep}{3pt}%
  \setlength{\parsep}{0pt}%
  \trivlist
  \item[\hskip\labelsep\bfseries\color{ScreenshotBlue}#1.]\ignorespaces
}{%
  \ifqedplaced\else\hfill$\square$\fi\endtrivlist
}

\newcommand{\qedhere}{\hfill$\square$\global\qedplacedtrue}

\usepackage{ stmaryrd } % for [[ and ]] symbols

\usepackage[noEnd=false,beginComment={\protect$\protect\smalltriangleright\protect$~}]{algpseudocodex} % for pseudocode algorithm typesetting
\tikzset{algpxIndentLine/.style={draw=gray}} 

\DeclareMathOperator{\card}{card}

% ==========================================
% DOCUMENT CONTENT
% ==========================================
\begin{document}

% --- SLIDE 1: TITLE SLIDE ---
\begin{frame}[plain]
    \centering
    \vfill
    {\fontsize{22}{28}\selectfont\bfseries\color{ScreenshotBlue} Multidimensional difference arrays \par}
    \vspace{0.1cm}
    {\fontsize{14}{18}\selectfont\color{TextBlack} via Discrete Calculus}
    \vfil
\end{frame}

% --- SLIDE 2: OVERVIEW ---
\begin{frame}{What is this talk about?}
    \begin{itemize}[<+->]
        \item A runtime optimization of \( L \) update queries on an array of size \( N \): 
            \begin{center}
                \Large
                \( \mathcal{O}(N\cdot L) \longrightarrow \mathcal{O}(N + L) \)
            \end{center}
        \item Simple heuristic for 1D, 2D and 3D arrays
        \item For general dimension: \emph{Discrete Calculus}
        \item Further optimizations
    \end{itemize}
\end{frame}

% --- SLIDE 3: 1D ARRAYS ---
\begin{frame}[fragile]{One-dimensional difference arrays}
    \begin{columns}[T]
	\begin{column}{0.45\textwidth}
	    \begin{tikzpicture}
		\def\cellWidth{0.6cm}    % Width of each array cell
		\def\cellHeight{0.6cm}   % Height of each array cell
		\def\indexGap{0.25cm}     % Vertical gap between indices and the array grid
		\def\rowBlockGap{1.9cm}  % Total vertical spacing from the top of one array row block to the next
		
		% ==========================================
		% Macro 1: Draw an array row
		% #1: label text (e.g., A:)
		% #2: comma-separated list of value/highlight pairs (e.g., 2/0, 6/1, 2/1) 
		%     where 0 = normal cell, 1 = green highlighted cell
		% #3: maximum index (e.g., 8 for 9 elements)
		% #4: absolute y-coordinate for the *top* of the array grid
		% ==========================================
		\newcommand{\drawArrayRow}[4]{
		    \begin{scope}[shift={(0,#4)}]
			% Label "A:" at the left
			\node at (-0.3, -\cellHeight/2) {{\ttfamily #1}};
			
			% Draw column indices above the array
			\foreach \i in {0,...,#3} {
			    \node[font=\footnotesize] at (\i*\cellWidth + \cellWidth/2, \indexGap) {\i};
			}
			
			% 1. Fill the backgrounds for highlighted cells first
			\foreach \val/\hl [count=\i from 0] in {#2} {
			    \ifnum\hl=1
				\fill[Green, opacity=0.1] (\i*\cellWidth, 0) rectangle ++(\cellWidth, -\cellHeight);
			    \fi
			}
			
			% 2. Draw the base black grid (prevents overlapping thick lines)
			\draw[thick] (0, 0) grid[xstep=\cellWidth, ystep=\cellHeight] ++(#3*\cellWidth+\cellWidth, -\cellHeight);
			
			% 3. Redraw the green borders over the grid for highlighted cells
			\foreach \val/\hl [count=\i from 0] in {#2} {
			    \ifnum\hl=1
				\draw[Green, thick] (\i*\cellWidth, 0) rectangle ++(\cellWidth, -\cellHeight);
			    \fi
			}
			
			% 4. Add the text nodes into the cells
			\foreach \val/\hl [count=\i from 0] in {#2} {
			    \ifnum\hl=1
				\node[math node,text=Green] at (\i*\cellWidth + \cellWidth/2, -\cellHeight/2) {\val};
			    \else
				\node[math node,text=black] at (\i*\cellWidth + \cellWidth/2, -\cellHeight/2) {\val};
			    \fi
			}
		    \end{scope}
		}
		
		% ==========================================
		% Macro 2: Draw a query arrow and text
		% #1: query text (e.g., Apply Q1)
		% #2: absolute y-coordinate to center the query block vertically
		% ==========================================
		\newcommand{\drawQueryBlock}[2]{
		    \begin{scope}[shift={(0,#2)}]
			 % Vertical arrow
			 \draw[-Stealth, thick] (-0.3, 0.1) -- (-0.3, -0.65);
			 % Query text to the right
			 \node[anchor=north west, inner sep=0pt] at (0, 0) {#1};
		    \end{scope}
		}

		\def\maxIndex{8} 

		\drawArrayRow{A:}{\ast/0, \ast/0, \ast/0, \ast/0, \ast/0, \ast/0, \ast/0, \ast/0, \ast/0}{\maxIndex}{0}
		
		\drawQueryBlock{Apply Q1}{-\rowBlockGap/2}
		
		\drawArrayRow{A:}{\ast/0, +1/1, +1/1, +1/1, \ast/0, \ast/0, \ast/0, \ast/0, \ast/0}{\maxIndex}{-\rowBlockGap}
		
		\drawQueryBlock{Apply Q2}{-\rowBlockGap - \rowBlockGap/2}
		
		\drawArrayRow{A:}{\ast/0, \ast/0, +2/1, +2/1, \ast/0, \ast/0, \ast/0, \ast/0, \ast/0}{\maxIndex}{-2*\rowBlockGap}
		
		\drawQueryBlock{Apply Q3}{-2*\rowBlockGap - \rowBlockGap/2}
		
		\drawArrayRow{A:}{\ast/0, \ast/0, \ast/0, \ast/0, \ast/0, -1/1, -1/1, -1/1, -1/1}{\maxIndex}{-3*\rowBlockGap}
	    \end{tikzpicture}
	\end{column}

    \begin{column}{0.45\textwidth}
	    Consider a one-dimensional array \( \texttt{A} \) of size \( 8 \). 

	    Apply the queries
	    \begin{itemize}
		\item Q1: Update elements at indices in \( \{1,2,3\} \) by \( +1 \)
		\item Q2: Update elements at indices in \( \{2,3\} \) by \( +2 \)
		\item Q3: Update elements at indices in \( \{5,6,7,8\} \) by \( -1 \)
	    \end{itemize}
	\end{column}

    \end{columns}
\end{frame}

\begin{frame}{One-dimensional difference arrays}
    Use two column layout and describe informal naive algorithm so that the runtime analysis is clearer. 

    \textbf{Runtime analysis:}
    \begin{itemize}
	\item One-dimensional array of size \( N \)
        \item \( Q \) queries for each of which the update range is at most of size \( N \)
    \end{itemize}
    Time complexity: \( \mathcal{O}(Q\cdot N) \)

    \vspace*{0.4cm}
    \textbf{We can do MUCH better!}
\end{frame}

\begin{frame}[fragile]{One-dimensional difference arrays}
    \begin{columns}[T]
	\begin{column}{0.4\textwidth}
	    We could also apply the query 
	    \begin{itemize}
		\item Q1: Update elements at indices in \( \{1,2,3\} \) by \( +1 \)
	    \end{itemize}
	    differently: 

	    \begin{enumerate}
		\item Initialize an array \texttt{d} of size 8
		\item Add \( +1 \) at index \( 1 \)
		\item Add \( -1 \) at index \( 4 \)
		\item Take the cumulative sum over the entire array \texttt{d}, that is \( \texttt{d}[i] = \sum_{j=0}^i \texttt{d}[j] \)
		\item Add \texttt{d} to \texttt{A} 
	    \end{enumerate}
	\end{column}

	\begin{column}{0.55\textwidth}
	    \centering
	    \begin{tikzpicture}
		\def\cellWidth{0.6cm}    % Width of each array cell
		\def\cellHeight{0.6cm}   % Height of each array cell
		\def\indexGap{0.25cm}     % Vertical gap between indices and the array grid
		\def\rowBlockGap{1.9cm}  % Total vertical spacing from the top of one array row block to the next
		
		% ==========================================
		% Macro 1: Draw an array row
		% #1: label text (e.g., A:)
		% #2: comma-separated list of value/highlight pairs (e.g., 2/0, 6/1, 2/1) 
		%     where 0 = normal cell, 1 = green highlighted cell
		% #3: maximum index (e.g., 8 for 9 elements)
		% #4: absolute y-coordinate for the *top* of the array grid
		% ==========================================
		\newcommand{\drawArrayRow}[4]{
		    \begin{scope}[shift={(0,#4)}]
			% Label "A:" at the left
			\node at (-0.3, -\cellHeight/2) {{\ttfamily #1}};
			
			% Draw column indices above the array
			\foreach \i in {0,...,#3} {
			    \node[font=\footnotesize] at (\i*\cellWidth + \cellWidth/2, \indexGap) {\i};
			}
			
			% 1. Fill the backgrounds for highlighted cells first
			\foreach \val/\hl [count=\i from 0] in {#2} {
			    \ifnum\hl=1
				\fill[Green, opacity=0.1] (\i*\cellWidth, 0) rectangle ++(\cellWidth, -\cellHeight);
			    \fi
			}
			
			% 2. Draw the base black grid (prevents overlapping thick lines)
			\draw[thick] (0, 0) grid[xstep=\cellWidth, ystep=\cellHeight] ++(#3*\cellWidth+\cellWidth, -\cellHeight);
			
			% 3. Redraw the green borders over the grid for highlighted cells
			\foreach \val/\hl [count=\i from 0] in {#2} {
			    \ifnum\hl=1
				\draw[Green, thick] (\i*\cellWidth, 0) rectangle ++(\cellWidth, -\cellHeight);
			    \fi
			}
			
			% 4. Add the text nodes into the cells
			\foreach \val/\hl [count=\i from 0] in {#2} {
			    \ifnum\hl=1
				\node[math node,text=Green] at (\i*\cellWidth + \cellWidth/2, -\cellHeight/2) {\val};
			    \else
				\node[math node,text=black] at (\i*\cellWidth + \cellWidth/2, -\cellHeight/2) {\val};
			    \fi
			}
		    \end{scope}
		}
		
		% ==========================================
		% Macro 2: Draw a query arrow and text
		% #1: query text (e.g., Apply Q1)
		% #2: absolute y-coordinate to center the query block vertically
		% ==========================================
		\newcommand{\drawQueryBlock}[2]{
		    \begin{scope}[shift={(0,#2)}]
			 % Vertical arrow
			 \draw[-Stealth, thick] (-0.3, 0.1) -- (-0.3, -0.65);
			 % Query text to the right
			 \node[anchor=north west, inner sep=0pt] at (0, 0) {#1};
		    \end{scope}
		}

		\def\maxIndex{8} 

		\drawArrayRow{D:}{0/0, 0/0, 0/0, 0/0, 0/0, 0/0, 0/0, 0/0, 0/0}{\maxIndex}{0}
		
		\drawQueryBlock{Mark Q1}{-\rowBlockGap/2}
		
		\drawArrayRow{D:}{0/0, +1/1, 0/0, 0/0, -1/1, 0/0, 0/0, 0/0, 0/0}{\maxIndex}{-\rowBlockGap}
		
		\drawQueryBlock{Cumulative Sum}{-\rowBlockGap - \rowBlockGap/2}
		
		\drawArrayRow{D:}{0/0, 1/0, 1/0, 1/0, 0/0, 0/0, 0/0, 0/0, 0/0}{\maxIndex}{-2*\rowBlockGap}
	    \end{tikzpicture}
	\end{column}
    \end{columns}
 
\end{frame}

\begin{frame}[fragile]{One-dimensional difference arrays}
    \begin{columns}[T]
	\begin{column}{0.45\textwidth}
	    This concept can be applied to multiple queries. 
	    \begin{itemize}
		\item Q1: Update elements at indices in \( \{1,2,3\} \) by \( +1 \)
		\item Q2: Update elements at indices in \( \{2,3\} \) by \( +2 \)
		\item Q3: Update elements at indices in \( \{5,6,7,8\} \) by \( -1 \)
	    \end{itemize}
	    Then form \( \texttt{d} = \texttt{d}_1 + \texttt{d}_2 + \texttt{d}_3 \). 
	\end{column}
	\begin{column}{0.5\textwidth}
	    \centering
	    \begin{tikzpicture}
		\def\cellWidth{0.6cm}    % Width of each array cell
		\def\cellHeight{0.6cm}   % Height of each array cell
		\def\indexGap{0.25cm}     % Vertical gap between indices and the array grid
		\def\rowBlockGap{1.2cm}  % Total vertical spacing from the top of one array row block to the next
		
		% ==========================================
		% Macro 1: Draw an array row
		% #1: label text (e.g., A:)
		% #2: comma-separated list of value/highlight pairs (e.g., 2/0, 6/1, 2/1) 
		%     where 0 = normal cell, 1 = green highlighted cell
		% #3: maximum index (e.g., 8 for 9 elements)
		% #4: absolute y-coordinate for the *top* of the array grid
		% ==========================================
		\newcommand{\drawArrayRow}[4]{
		    \begin{scope}[shift={(0,#4)}]
			% Label "A:" at the left
			\node at (-0.3, -\cellHeight/2) {#1};
			
			% Draw column indices above the array
			\foreach \i in {0,...,#3} {
			    \node[font=\footnotesize] at (\i*\cellWidth + \cellWidth/2, \indexGap) {\i};
			}
			
			% 1. Fill the backgrounds for highlighted cells first
			\foreach \val/\hl [count=\i from 0] in {#2} {
			    \ifnum\hl=1
				\fill[Green, opacity=0.1] (\i*\cellWidth, 0) rectangle ++(\cellWidth, -\cellHeight);
			    \fi
			}
			
			% 2. Draw the base black grid (prevents overlapping thick lines)
			\draw[thick] (0, 0) grid[xstep=\cellWidth, ystep=\cellHeight] ++(#3*\cellWidth+\cellWidth, -\cellHeight);
			
			% 3. Redraw the green borders over the grid for highlighted cells
			\foreach \val/\hl [count=\i from 0] in {#2} {
			    \ifnum\hl=1
				\draw[Green, thick] (\i*\cellWidth, 0) rectangle ++(\cellWidth, -\cellHeight);
			    \fi
			}
			
			% 4. Add the text nodes into the cells
			\foreach \val/\hl [count=\i from 0] in {#2} {
			    \ifnum\hl=1
				\node[math node,text=Green] at (\i*\cellWidth + \cellWidth/2, -\cellHeight/2) {\val};
			    \else
				\node[math node,text=black] at (\i*\cellWidth + \cellWidth/2, -\cellHeight/2) {\val};
			    \fi
			}
		    \end{scope}
		}
		
		% ==========================================
		% Macro 2: Draw a query arrow and text
		% #1: query text (e.g., Apply Q1)
		% #2: absolute y-coordinate to center the query block vertically
		% ==========================================
		\newcommand{\drawQueryBlock}[2]{
		    \begin{scope}[shift={(0,#2)}]
			 % Vertical arrow
			 \draw[-Stealth, thick] (-0.3, 0.1) -- (-0.3, -0.65);
			 % Query text to the right
			 \node[anchor=north west, inner sep=0pt] at (0, 0) {#1};
		    \end{scope}
		}

		\def\maxIndex{8} 

		\drawArrayRow{$\texttt{d}_1$:}{0/0, +1/0, 0/0, 0/0, -1/0, 0/0, 0/0, 0/0, 0/0}{\maxIndex}{0}
		
		\drawArrayRow{$\texttt{d}_2$:}{0/0, 0/0, +2/0, 0/0, -2/0, 0/0, 0/0, 0/0, 0/0}{\maxIndex}{-\rowBlockGap}
		
		\drawArrayRow{$\texttt{d}_3$:}{0/0, 0/0, 0/0, 0/0, 0/0, -1/0, 0/0, 0/0, 0/0}{\maxIndex}{-2*\rowBlockGap}

		\drawQueryBlock{$+$}{-2.75*\rowBlockGap}

		\drawArrayRow{$\texttt{d}$:}{0/0, +1/0, +2/0, 0/0, -3/0, -1/0, 0/0, 0/0, 0/0}{\maxIndex}{-3.5*\rowBlockGap}

	    \end{tikzpicture}
	\end{column}
    \end{columns}
\end{frame}

\begin{frame}[fragile]{One-dimensional difference arrays}
    \begin{columns}[T]
	\begin{column}{0.45\textwidth}
	    The cumulative sum distributes over \( \texttt{d} = \texttt{d}_1 + \texttt{d}_2 + \texttt{d}_3 \):
	    \begin{align*}
		&\sum_{j=0}^i \texttt{d}[j] \\
		&= \sum_{j=0}^i \left(\texttt{d}_1[j] + \texttt{d}_3[j] + \texttt{d}_3[j] \right) \\
		&= \sum_{j=0}^i \texttt{d}_1[j] + \sum_{j=0}^i \texttt{d}_2[j] + \sum_{j=0}^i \texttt{d}_3[j]
	    \end{align*}
	    Hence, \( \Sigma \texttt{d} = \Sigma\texttt{d}_1 + \Sigma\texttt{d}_2 + \Sigma\texttt{d}_3 \). 

	    Finally add \( \Sigma\texttt{d} \) to \texttt{A}. 
	\end{column}
	\begin{column}{0.5\textwidth}
	    \centering
	    \begin{tikzpicture}
		\def\cellWidth{0.6cm}    % Width of each array cell
		\def\cellHeight{0.6cm}   % Height of each array cell
		\def\indexGap{0.25cm}     % Vertical gap between indices and the array grid
		\def\rowBlockGap{1.2cm}  % Total vertical spacing from the top of one array row block to the next
		
		% ==========================================
		% Macro 1: Draw an array row
		% #1: label text (e.g., A:)
		% #2: comma-separated list of value/highlight pairs (e.g., 2/0, 6/1, 2/1) 
		%     where 0 = normal cell, 1 = green highlighted cell
		% #3: maximum index (e.g., 8 for 9 elements)
		% #4: absolute y-coordinate for the *top* of the array grid
		% ==========================================
		\newcommand{\drawArrayRow}[4]{
		    \begin{scope}[shift={(0,#4)}]
			% Label "A:" at the left
			\node[anchor=east] at (0, -\cellHeight/2) {#1};
			
			% Draw column indices above the array
			\foreach \i in {0,...,#3} {
			    \node[font=\footnotesize] at (\i*\cellWidth + \cellWidth/2, \indexGap) {\i};
			}
			
			% 1. Fill the backgrounds for highlighted cells first
			\foreach \val/\hl [count=\i from 0] in {#2} {
			    \ifnum\hl=1
				\fill[Green, opacity=0.1] (\i*\cellWidth, 0) rectangle ++(\cellWidth, -\cellHeight);
			    \fi
			}
			
			% 2. Draw the base black grid (prevents overlapping thick lines)
			\draw[thick] (0, 0) grid[xstep=\cellWidth, ystep=\cellHeight] ++(#3*\cellWidth+\cellWidth, -\cellHeight);
			
			% 3. Redraw the green borders over the grid for highlighted cells
			\foreach \val/\hl [count=\i from 0] in {#2} {
			    \ifnum\hl=1
				\draw[Green, thick] (\i*\cellWidth, 0) rectangle ++(\cellWidth, -\cellHeight);
			    \fi
			}
			
			% 4. Add the text nodes into the cells
			\foreach \val/\hl [count=\i from 0] in {#2} {
			    \ifnum\hl=1
				\node[math node,text=Green] at (\i*\cellWidth + \cellWidth/2, -\cellHeight/2) {\val};
			    \else
				\node[math node,text=black] at (\i*\cellWidth + \cellWidth/2, -\cellHeight/2) {\val};
			    \fi
			}
		    \end{scope}
		}
		
		% ==========================================
		% Macro 2: Draw a query arrow and text
		% #1: query text (e.g., Apply Q1)
		% #2: absolute y-coordinate to center the query block vertically
		% ==========================================
		\newcommand{\drawQueryBlock}[2]{
		    \begin{scope}[shift={(0,#2)}]
			 % Vertical arrow
			 \draw[-Stealth, thick] (-0.3, 0.1) -- (-0.3, -0.65);
			 % Query text to the right
			 \node[anchor=north west, inner sep=0pt] at (0, 0) {#1};
		    \end{scope}
		}

		\def\maxIndex{8} 

		\drawArrayRow{$\Sigma\texttt{d}_1$:}{0/0, 1/0, 1/0, 1/0, 0/0, 0/0, 0/0, 0/0, 0/0}{\maxIndex}{0}
		
		\drawArrayRow{$\Sigma\texttt{d}_2$:}{0/0, 0/0, 2/0, 2/0, 0/0, 0/0, 0/0, 0/0, 0/0}{\maxIndex}{-\rowBlockGap}
		
		\drawArrayRow{$\Sigma\texttt{d}_3$:}{0/0, 0/0, 0/0, 0/0, 0/0, -1/0, -1/0, -1/0, -1/0}{\maxIndex}{-2*\rowBlockGap}

		\drawQueryBlock{$+$}{-2.75*\rowBlockGap}

		\drawArrayRow{$\Sigma\texttt{d}$:}{0/0, 1/0, 3/0, 3/0, 0/0, -1/0, -1/0, -1/0, -1/0}{\maxIndex}{-3.5*\rowBlockGap}

	    \end{tikzpicture}
	\end{column}
    \end{columns}
\end{frame}

\begin{frame}{One-dimensional difference arrays}
    \begin{columns}[T]
	\begin{column}{0.5\textwidth}
	    \textbf{Non-formal difference array algorithm:}

	    Input: Array \texttt{A} of size \( N \) and \( L \) queries of size at most \( N \)
	    \begin{enumerate}
		\item Initialize array \texttt{d} of size \( N \) 
		\item For every query add two markers to \texttt{d}
		\item Form the cumulative sum \( \Sigma\texttt{d} \)
		\item Add \( \Sigma\texttt{d} \) to \texttt{A}
	    \end{enumerate}
	\end{column}
	\begin{column}{0.45\textwidth}
	    \textbf{Runtime analysis:}

	    \begin{itemize}
		\item Initializing \texttt{d} in \( \mathcal{O}(N) \)
		\item Per query placing two markers, i.e. in total \( \mathcal{O}(L) \) markers
		\item Forming \( \Sigma\texttt{d} \) in \( \mathcal{O}(N) \) by computing 
		    \begin{equation*}
			\sum_{j=0}^i \texttt{d}[j] = \texttt{d}[i] + \sum_{j=0}^{i-1} \texttt{d}[j]
		    \end{equation*}
		    (\(N-2\) additions, \(N-1\) reads)
		\item Adding \( \Sigma\texttt{d} \) to \texttt{A} in \( \mathcal{O}(N) \). 
	    \end{itemize}
	    In total: \( \mathcal{O}(N + L) \). 
	\end{column}
    \end{columns}
\end{frame}

\begin{frame}{Two-dimensional difference arrays}
    This concept can easily be generalized to 2D arrays. 

    \begin{columns}[T]
	\begin{column}{0.5\textwidth}
	    \begin{tikzpicture}[
		% Style for the grid lines
		gridline/.style={thick, draw=black},
		% Style for the highlighted 3x5 green box
		greenbox/.style={draw=green!60!black, ultra thick, fill=green!60!black, fill opacity=0.1},
		% Style for the asterisks
		asteriskstyle/.style={font=\large, inner sep=0pt},
		scale=0.75
	    ]

		% 1. Draw the 5x8 Grid
		% 8 columns wide (0 to 8), 5 rows high (0 to 5)
		\draw[gridline] (0,0) grid (8,5);
		
		% 2. Place an asterisk (*) in the center of each of the 40 boxes
		% Loop through columns (x) from 1 to 8, and rows (y) from 1 to 5
		\foreach \x in {1,...,8} {
		    \node[asteriskstyle] at (\x-0.5, 0.5) {$\ast$};
		}
		\foreach \x in {1,...,8} {
		    \node[asteriskstyle] at (\x-0.5, 4.5) {$\ast$};
		}
		\foreach \x in {1,7,8} {
		    \foreach \y in {1,...,5} {
			\node[asteriskstyle] at (\x-0.5, \y-0.5) {$\ast$};
		    }
		}
		\foreach \x in {2,...,6} {
		    \foreach \y in {2,...,4} {
			\node[asteriskstyle] at (\x-0.5,\y-0.5) {$+C$};
		    }
		}
		
		% 3. Draw the green box around a 3x5 area
		% This highlights the area spanning columns 2 to 6 (width 5) 
		% and rows 2 to 4 (height 3) as an example. 
		\draw[greenbox] (1,1) rectangle (6,4);
	    \end{tikzpicture}

	    Naively applying the constant update to a \( 5 \times 8 \) array \texttt{A}.
	\end{column}
	\begin{column}{0.45\textwidth}
	    \begin{tikzpicture}[
		% Style for the grid lines
		gridline/.style={thick, draw=black},
		% Style for the highlighted 3x5 green box
		greenbox/.style={draw=green!60!black, ultra thick, fill=green!60!black, fill opacity=0.1},
		% Style for the asterisks
		asteriskstyle/.style={font=\large, inner sep=0pt},
		scale=0.75
	    ]

		% 1. Draw the 5x8 Grid
		% 8 columns wide (0 to 8), 5 rows high (0 to 5)
		\draw[gridline] (0,0) grid (8,5);
		
		% 3. Draw the green box around a 3x5 area
		% This highlights the area spanning columns 2 to 6 (width 5) 
		% and rows 2 to 4 (height 3) as an example. 
		\draw[greenbox] (1,1) rectangle (6,4);

		\node[asteriskstyle] at (1.5, 1.5) {$+C$};
		\node[asteriskstyle] at (6.5, 1.5) {$-C$};
		\node[asteriskstyle] at (6.5, 4.5) {$+C$};
		\node[asteriskstyle] at (1.5, 4.5) {$-C$};
	    \end{tikzpicture}

	    Placing \emph{four} markers in the difference array. 
	\end{column}
    \end{columns}
\end{frame}

\begin{frame}{Two-dimensional difference arrays}
    Taking cumulative sums in both directions yields the correct upates.

    \begin{columns}[T]
	\begin{column}{0.5\textwidth}
	    \vfill

	    \begin{tikzpicture}[
		% Style for the grid lines
		gridline/.style={thick, draw=black},
		% Style for the highlighted 3x5 green box
		greenbox/.style={draw=green!60!black, ultra thick, fill=green!60!black, fill opacity=0.1},
		% Style for the asterisks
		asteriskstyle/.style={font=\large, inner sep=0pt},
		scale=0.75
	    ]

		% 1. Draw the 5x8 Grid
		% 8 columns wide (0 to 8), 5 rows high (0 to 5)
		\draw[gridline] (0,0) grid (8,5);
		
		% 3. Draw the green box around a 3x5 area
		% This highlights the area spanning columns 2 to 6 (width 5) 
		% and rows 2 to 4 (height 3) as an example. 
		\draw[greenbox] (1,1) rectangle (6,4);

		\node[asteriskstyle] at (1.5, 1.5) {$+C$};
		\node[asteriskstyle] at (1.5, 2.5) {$+C$};
		\node[asteriskstyle] at (1.5, 3.5) {$+C$};
		\node[asteriskstyle] at (6.5, 1.5) {$-C$};
		\node[asteriskstyle] at (6.5, 2.5) {$-C$};
		\node[asteriskstyle] at (6.5, 3.5) {$-C$};

		\draw[-Stealth,thick] (-0.5,0) -- (-0.5,5);
		\draw[-Stealth,thick,white] (0,-0.5) -- (8,-0.5);
	    \end{tikzpicture}

	    \vfill
	\end{column}
	\begin{column}{0.45\textwidth}
	    \vfill
	    \begin{tikzpicture}[
		% Style for the grid lines
		gridline/.style={thick, draw=black},
		% Style for the highlighted 3x5 green box
		greenbox/.style={draw=green!60!black, ultra thick, fill=green!60!black, fill opacity=0.1},
		% Style for the asterisks
		asteriskstyle/.style={font=\large, inner sep=0pt},
		scale=0.75
	    ]

		% 1. Draw the 5x8 Grid
		% 8 columns wide (0 to 8), 5 rows high (0 to 5)
		\draw[gridline] (0,0) grid (8,5);
		
		% 3. Draw the green box around a 3x5 area
		% This highlights the area spanning columns 2 to 6 (width 5) 
		% and rows 2 to 4 (height 3) as an example. 
		\draw[greenbox] (1,1) rectangle (6,4);

		\foreach \x in {2,...,6} {
		    \foreach \y in {2,...,4} {
			\node[asteriskstyle] at (\x-0.5,\y-0.5) {$+C$};
		    }
		}

		\draw[-Stealth,thick] (0,-0.5) -- (8,-0.5);
	    \end{tikzpicture}
	    \vfill
	\end{column}
    \end{columns}
\end{frame}

\begin{frame}{\(d\)-dimensional difference arrays?}
    \begin{itemize}
	\item This heuristic can easily be extended to \( d \)-dimensional arrays by placing \( 2^d \) markers for every query. 

	\item However for dimensions \( d \geq 4 \) this is hard to conceptially verify. 

	\item Furthermore, "naive" proofs without any additional tools are bound to get very messy. 

	\item This is where \emph{Discrete Calculus} comes into play. 

	\item To leaverage these tools, we first need to formally describe our situation.
    \end{itemize}

\end{frame}

\begin{frame}{The formal framework}
    We view arrays as maps over \emph{finite grid domains}. 

    \begin{definition*}{}
	A finite subset \( D \subseteq \mathbb{Z}^d \) is called \emph{finite grid domain} if there exist \( n_1,\dots,n_d \in \mathbb{Z} \) and \( m_1, \dots, m_d \in \mathbb{Z} \) with \( n_i \leq m_i \) for all \( i \in \{1, \dots, d\} \) such that \( D = \prod_{i=1}^d \llbracket n_i, m_i \rrbracket \).
	A \emph{discrete map} is a map \( f : D \to \mathbb{R} \) over a finite grid domain \( D \subseteq \mathbb{Z}^d \). 

	A \emph{\(d\)-dimensional array} is a discrete function \( f : D \to \mathbb R \) for which the finite grid domain \( D \) is of the form \( D = \prod_{i=1}^d \llbracket 0, m_i \rrbracket \subseteq \mathbb Z^d \) where \( m_i \in \mathbb N \) for all \( 1 \leq i \leq d \).  
    \end{definition*}

    The notation \( \llbracket a, b \rrbracket \) for \( a\leq b \in \mathbb{Z} \) means the \emph{discrete interval}
    \begin{equation*}
	\llbracket a, b \rrbracket := \{ m \in \mathbb{Z} \mid a \leq m \leq b \}. 
    \end{equation*}
\end{frame}

\begin{frame}
    \begin{definition*}{}
	Let \( D = \prod_{i=1}^d \llbracket n_i, m_i \rrbracket \subseteq \mathbb{Z}^d \) be a finite grid domain. 
	A \emph{constant update query} \( q \) on \( D \) is a discrete function \( q : D \to \mathbb R \) such that there exists a subset \( D_q = \prod_{i=1}^d \llbracket a_i, b_i \rrbracket \subseteq D \) called the \emph{update range} of \( q \) and a constant \( C \in \mathbb R \) called the \emph{update value} of \( q \) with 
	\begin{equation*}
	    q(k_1, \dots, k_d) := \begin{cases}
		C, & \text{if } (k_1,\dots,k_d) \in D_q, \\
		0, & \text{else}
	    \end{cases}
	\end{equation*}
	for \( (k_1,\dots,k_d) \in D \), that is \( q = C \cdot \mathbb{1}_{D_q} \). 
	The \emph{size} of a constant update query \( q : D \to \mathbb R \) is the cardinality of the update range \( D_q \). 
    \end{definition*}
\end{frame}

\begin{frame}
    Let \( f : D \to \mathbb R \) with \( D = \prod_{i=1}^d \llbracket n_i, m_i \rrbracket \) be a discrete function and \( q_1, \dots, q_L \) a series of constant update queries. 

    \vspace{.3cm}

    \emph{Applying the series of queries to \( f \)} means forming the updated discrete function \( F : D \to \mathbb R \) as 
    \begin{equation*}
	F := f + \sum_{i=1}^L q_i.
    \end{equation*}
\end{frame}

\begin{frame}{A naive algorithm}
    \begin{algorithmic}[1]
	\Statex\hspace{-2em} \textbf{Input:} A discrete map \( f : D \to \mathbb R \) with \( D = \prod_{i=1}^d \llbracket n_i, m_i \rrbracket \) and a series of constant update queries \( q_1, \dots, q_L \) on \( f \) over the update ranges \( D_j = \prod_{i=1}^d \llbracket a_i^{(j)}, b_i^{(j)} \rrbracket \) with update values \( C_j \) for \( 1 \leq j \leq L \)

	\Statex\hspace{-2em} \textbf{Output:} The updated function \( F = f + \sum_{i=1}^L q_i \)

	\Statex

	\State \( F \gets \) The map \( f : D \to \mathbb{R} \) \hspace{2em}
	\Comment{For an in situ implementation set \( F = f \) without any copy operations}

	\For{\( j \in \llbracket 1, L \rrbracket \)}
	    \For{\( k \in D_j \)}
		\State \( F(k) \gets F(k) + C_j \)
	    \EndFor
	\EndFor

	\Return \( F \)
    \end{algorithmic}
\end{frame}

\begin{frame}
    \vspace{.5cm}
    \begin{tikzpicture}
	\draw[help lines, color=gray!30, step=1cm] (0,0) grid (12, 7);

	% 2. Draw Axes
	\draw[->, thick] (-0.15, 0) -- (12, 0); % X-axis
	\draw[->, thick] (0, -0.15) -- (0, 7);  % Y-axis

	% 3. Draw Axis Ticks and Labels
	% X-axis ticks (at integer coordinates)
	\draw[thick] (2, -0.1) -- (2, 0.1) node[below=5pt] {$a_1^{(1)}$};
	\draw[thick] (6, -0.1) -- (6, 0.1) node[below=5pt] {$b_1^{(1)}$};

	\draw[thick] (4, -0.1) -- (4, 0.1) node[below=5pt] {$a_1^{(2)}$};
	\draw[thick] (7, -0.1) -- (7, 0.1) node[below=5pt] {$b_1^{(2)}$};

	\draw[thick] (9, -0.1) -- (9, 0.1) node[below=5pt] {$a_1^{(3)}$};
	\draw[thick] (10, -0.1) -- (10, 0.1) node[below=5pt] {$b_1^{(3)}$};

	% Y-axis ticks (at integer coordinates)
	\draw[thick] (-0.1, 3) -- (0.1, 3) node[left=5pt] {$a_2^{(1)}$};
	\draw[thick] (-0.1, 5) -- (0.1, 5) node[left=5pt] {$b_2^{(1)}$};

	\draw[thick] (-0.1, 2) -- (0.1, 2) node[left=5pt] {$a_2^{(2)}$};
	\draw[thick] (-0.1, 4) -- (0.1, 4) node[left=5pt] {$b_2^{(2)}$};

	% 4. Draw Rectangles (Edges exactly between grid points)
	% y-span is from 1.5 to 6.5
	\fill[gray, opacity=0.2] (2-0.5, 3-0.5) rectangle (6+0.5, 6-0.5);  % q1 bounding box
	\fill[gray, opacity=0.3] (4-0.5, 2-0.5) rectangle (7+0.5, 5-0.5);  % q1 bounding box
	\fill[gray, opacity=0.4] (9-0.5, 3-0.5) rectangle (10+0.5, 6-0.5);  % q1 bounding box

	% 5. Add Rectangle Area Labels
	\node at (2.5, 2) {$q_1$};
	\node at (5.5, 1) {$q_2$};
	\node at (9.5, 2) {$q_3$};
	
	% q1 updates
	\node[anchor=center, inner sep=2pt] at (2,3) {$+C_1$};
	\node[anchor=center, inner sep=2pt] at (2,4) {$+C_1$};
	\node[anchor=center, inner sep=2pt] at (2,5) {$+C_1$};
	\node[anchor=center, inner sep=2pt] at (3,3) {$+C_1$};
	\node[anchor=center, inner sep=2pt] at (3,4) {$+C_1$};
	\node[anchor=center, inner sep=2pt] at (3,5) {$+C_1$};
	\node[anchor=center, inner sep=2pt] at (4,5) {$+C_1$};
	\node[anchor=center, inner sep=2pt] at (5,5) {$+C_1$};
	\node[anchor=center, inner sep=2pt] at (6,5) {$+C_1$};

	% q2 updates
	\node[anchor=center, inner sep=2pt] at (4,2) {$+C_2$};
	\node[anchor=center, inner sep=2pt] at (5,2) {$+C_2$};
	\node[anchor=center, inner sep=2pt] at (6,2) {$+C_2$};
	\node[anchor=center, inner sep=2pt] at (7,2) {$+C_2$};
	\node[anchor=center, inner sep=2pt] at (7,3) {$+C_2$};
	\node[anchor=center, inner sep=2pt] at (7,4) {$+C_2$};

	% q1/q2 overlap
	\node[anchor=center, align=left, inner sep=2pt] at (4,3) {$+C_1$\\[-0.5ex]$+C_2$};
	\node[anchor=center, align=left, inner sep=2pt] at (5,3) {$+C_1$\\[-0.5ex]$+C_2$};
	\node[anchor=center, align=left, inner sep=2pt] at (6,3) {$+C_1$\\[-0.5ex]$+C_2$};
	\node[anchor=center, align=left, inner sep=2pt] at (4,4) {$+C_1$\\[-0.5ex]$+C_2$};
	\node[anchor=center, align=left, inner sep=2pt] at (5,4) {$+C_1$\\[-0.5ex]$+C_2$};
	\node[anchor=center, align=left, inner sep=2pt] at (6,4) {$+C_1$\\[-0.5ex]$+C_2$};

	% q3 updates
	\node[anchor=center, inner sep=2pt] at (9,3) {$+C_3$};
	\node[anchor=center, inner sep=2pt] at (9,4) {$+C_3$};
	\node[anchor=center, inner sep=2pt] at (9,5) {$+C_3$};
	\node[anchor=center, inner sep=2pt] at (10,3) {$+C_3$};
	\node[anchor=center, inner sep=2pt] at (10,4) {$+C_3$};
	\node[anchor=center, inner sep=2pt] at (10,5) {$+C_3$};
    \end{tikzpicture}
\end{frame}

\begin{frame}
    \begin{proposition*}{}
	Let \( f : D \subseteq \mathbb{Z}^d \to \mathbb{R} \) be a discrete function with \( D = \prod_{i=1}^d \llbracket n_i, m_i \rrbracket \) and \( N = \card(D) \), and let \( q_1, \dots, q_Q \) be a series of \( L \) constant update queries on \( f \) each of size at most \( M \). 

	\begin{itemize}
	    \item Assuming that for each \( k \in D \) the function value \( f(k) \) can be obtained in constant time, the naive algorithm computes the updated function \( F = f + \sum_{i=1}^Q q_i \) in \( \mathcal{O}(N + LM) \). 

	    \item If the algorithm can be implemented in situ, its runtime lies in \( \mathcal{O}(LM) \). 
	\end{itemize}
    \end{proposition*}
\end{frame}

\begin{frame}
    \vspace{.5cm}
    \begin{tikzpicture}
	\draw[help lines, color=gray!30, step=1cm] (0,0) grid (12, 7);

	% 2. Draw Axes
	\draw[->, thick] (-0.15, 0) -- (12, 0); % X-axis
	\draw[->, thick] (0, -0.15) -- (0, 7);  % Y-axis

	% 3. Draw Axis Ticks and Labels
	% X-axis ticks (at integer coordinates)
	\draw[thick] (2, -0.1) -- (2, 0.1) node[below=5pt] {$a_1^{(1)}$};
	\draw[thick] (6, -0.1) -- (6, 0.1) node[below=5pt] {$b_1^{(1)}$};

	\draw[thick] (4, -0.1) -- (4, 0.1) node[below=5pt] {$a_1^{(2)}$};
	\draw[thick] (7, -0.1) -- (7, 0.1) node[below=5pt] {$b_1^{(2)}$};

	\draw[thick] (9, -0.1) -- (9, 0.1) node[below=5pt] {$a_1^{(3)}$};
	\draw[thick] (10, -0.1) -- (10, 0.1) node[below=5pt] {$b_1^{(3)}$};

	% Y-axis ticks (at integer coordinates)
	\draw[thick] (-0.1, 3) -- (0.1, 3) node[left=5pt] {$a_2^{(1)}$};
	\draw[thick] (-0.1, 5) -- (0.1, 5) node[left=5pt] {$b_2^{(1)}$};

	\draw[thick] (-0.1, 2) -- (0.1, 2) node[left=5pt] {$a_2^{(2)}$};
	\draw[thick] (-0.1, 4) -- (0.1, 4) node[left=5pt] {$b_2^{(2)}$};

	% 4. Draw Rectangles (Edges exactly between grid points)
	% y-span is from 1.5 to 6.5
	\fill[gray, opacity=0.2] (2-0.5, 3-0.5) rectangle (6+0.5, 6-0.5);  % q1 bounding box
	\fill[gray, opacity=0.3] (4-0.5, 2-0.5) rectangle (7+0.5, 5-0.5);  % q1 bounding box
	\fill[gray, opacity=0.4] (10-0.5, 3-0.5) rectangle (12+0.5, 6-0.5);  % q1 bounding box

	% 5. Add Rectangle Area Labels
	\node at (2.5, 2) {$q_1$};
	\node at (5.5, 1) {$q_2$};
	\node at (11, 2) {$q_3$};
	
	% q1 updates
	\node[anchor=center, inner sep=2pt] at (2,3) {$+C_1$};
	\node[anchor=center, inner sep=2pt] at (2,4) {$+C_1$};
	\node[anchor=center, inner sep=2pt] at (2,5) {$+C_1$};
	\node[anchor=center, inner sep=2pt] at (3,3) {$+C_1$};
	\node[anchor=center, inner sep=2pt] at (3,4) {$+C_1$};
	\node[anchor=center, inner sep=2pt] at (3,5) {$+C_1$};
	\node[anchor=center, inner sep=2pt] at (4,5) {$+C_1$};
	\node[anchor=center, inner sep=2pt] at (5,5) {$+C_1$};
	\node[anchor=center, inner sep=2pt] at (6,5) {$+C_1$};

	% q2 updates
	\node[anchor=center, inner sep=2pt] at (4,2) {$+C_2$};
	\node[anchor=center, inner sep=2pt] at (5,2) {$+C_2$};
	\node[anchor=center, inner sep=2pt] at (6,2) {$+C_2$};
	\node[anchor=center, inner sep=2pt] at (7,2) {$+C_2$};
	\node[anchor=center, inner sep=2pt] at (7,3) {$+C_2$};
	\node[anchor=center, inner sep=2pt] at (7,4) {$+C_2$};

	% q1/q2 overlap
	\node[anchor=center, align=left, inner sep=2pt] at (4,3) {$+C_1$\\[-0.5ex]$+C_2$};
	\node[anchor=center, align=left, inner sep=2pt] at (5,3) {$+C_1$\\[-0.5ex]$+C_2$};
	\node[anchor=center, align=left, inner sep=2pt] at (6,3) {$+C_1$\\[-0.5ex]$+C_2$};
	\node[anchor=center, align=left, inner sep=2pt] at (4,4) {$+C_1$\\[-0.5ex]$+C_2$};
	\node[anchor=center, align=left, inner sep=2pt] at (5,4) {$+C_1$\\[-0.5ex]$+C_2$};
	\node[anchor=center, align=left, inner sep=2pt] at (6,4) {$+C_1$\\[-0.5ex]$+C_2$};

	% q3 updates
	\node[anchor=center, inner sep=2pt] at (10,3) {$+C_3$};
	\node[anchor=center, inner sep=2pt] at (10,4) {$+C_3$};
	\node[anchor=center, inner sep=2pt] at (10,5) {$+C_3$};
	\node[anchor=center, inner sep=2pt] at (11,3) {$+C_3$};
	\node[anchor=center, inner sep=2pt] at (11,4) {$+C_3$};
	\node[anchor=center, inner sep=2pt] at (11,5) {$+C_3$};
	\node[anchor=center, inner sep=2pt] at (12,3) {$+C_3$};
	\node[anchor=center, inner sep=2pt] at (12,4) {$+C_3$};
	\node[anchor=center, inner sep=2pt] at (12,5) {$+C_3$};
    \end{tikzpicture}
\end{frame}

\begin{frame}
    \vspace{.5cm}
    \begin{tikzpicture}
	\draw[help lines, color=gray!30, step=1cm] (0,0) grid (12, 7);

	% 2. Draw Axes
	\draw[->, thick] (-0.15, 0) -- (12, 0); % X-axis
	\draw[->, thick] (0, -0.15) -- (0, 7);  % Y-axis

	% 3. Draw Axis Ticks and Labels
	% X-axis ticks (at integer coordinates)
	\draw[thick] (2, -0.1) -- (2, 0.1) node[below=5pt] {$a_1^{(1)}$};
	\draw[thick] (6, -0.1) -- (6, 0.1) node[below=5pt] {$b_1^{(1)}$};

	\draw[thick] (4, -0.1) -- (4, 0.1) node[below=5pt] {$a_1^{(2)}$};
	\draw[thick] (7, -0.1) -- (7, 0.1) node[below=5pt] {$b_1^{(2)}$};

	\draw[thick] (9, -0.1) -- (9, 0.1) node[below=5pt] {$a_1^{(3)}$};
	\draw[thick] (10, -0.1) -- (10, 0.1) node[below=5pt] {$b_1^{(3)}$};

	% Y-axis ticks (at integer coordinates)
	\draw[thick] (-0.1, 3) -- (0.1, 3) node[left=5pt] {$a_2^{(1)}$};
	\draw[thick] (-0.1, 5) -- (0.1, 5) node[left=5pt] {$b_2^{(1)}$};

	\draw[thick] (-0.1, 2) -- (0.1, 2) node[left=5pt] {$a_2^{(2)}$};
	\draw[thick] (-0.1, 4) -- (0.1, 4) node[left=5pt] {$b_2^{(2)}$};

	% 4. Draw Rectangles (Edges exactly between grid points)
	% y-span is from 1.5 to 6.5
	\fill[gray, opacity=0.2] (2-0.5, 3-0.5) rectangle (6+0.5, 6-0.5);  % q1 bounding box
	\fill[gray, opacity=0.3] (4-0.5, 2-0.5) rectangle (7+0.5, 5-0.5);  % q1 bounding box
	\fill[gray, opacity=0.4] (10-0.5, 3-0.5) rectangle (12+0.5, 6-0.5);  % q1 bounding box

	% 5. Add Rectangle Area Labels
	\node at (2.5, 2) {$q_1$};
	\node at (5.5, 1) {$q_2$};
	\node at (11, 2) {$q_3$};
	
	% q1 marker
	\node[anchor=center, inner sep=2pt] at (2,3) {$+C_1$};
	\node[anchor=center, inner sep=2pt] at (7,3) {$-C_1$};
	\node[anchor=center, inner sep=2pt] at (7,6) {$+C_1$};
	\node[anchor=center, inner sep=2pt] at (2,6) {$-C_1$};

	% q2 marker
	\node[anchor=center, inner sep=2pt] at (4,2) {$+C_2$};
	\node[anchor=center, inner sep=2pt] at (8,2) {$-C_2$};
	\node[anchor=center, inner sep=2pt] at (8,5) {$+C_2$};
	\node[anchor=center, inner sep=2pt] at (4,5) {$-C_2$};

	% q3 updates
	\node[anchor=center, inner sep=2pt] at (10,3) {$+C_3$};
	%\node[anchor=center, inner sep=2pt] at (11,3) {$-C_3$};
	\node[anchor=center, inner sep=2pt] at (10,6) {$-C_3$};
	%\node[anchor=center, inner sep=2pt] at (11,6) {$+C_3$};
    \end{tikzpicture}
\end{frame}

% \begin{frame}
%     \textbf{What are difference arrays formally?}
% 
%     \begin{itemize}
% 	\item Finite grid domain \( D = \llbracket n_1, m_1 \rrbracket \times \llbracket n_2, m_2 \rrbracket \)
% 	\item Constant update queries \( q_1, q_2 \) and \( q_3 \) with update ranges 
% 	    \begin{align*}
% 		D_1 &= \llbracket a_1^{(1)}, b_1^{(1)} \rrbracket \times \llbracket a_2^{(1)}, b_2^{(1)} \rrbracket, \\
% 		D_2 &= \llbracket a_1^{(2)}, b_1^{(2)} \rrbracket \times \llbracket a_2^{(2)}, b_2^{(2)} \rrbracket \\
% 		D_3 &= \llbracket a_1^{(3)}, b_1^{(3)} \rrbracket \times \llbracket a_2^{(3)}, b_2^{(3)} \rrbracket, 
% 	    \end{align*}
% 	    and update values \( C_1, C_2 \) and \( C_3 \)
% 	\item Form the difference array as 
% 	    \begin{align*}
% 		\Delta[q_1,q_2,q_3] = 
% 	    \end{align*}
%     \end{itemize}
% \end{frame}

\begin{frame}
    \vspace{.5cm}
    \begin{tikzpicture}
	\draw[help lines, color=gray!30, step=1cm] (0,0) grid (12, 7);

	% 2. Draw Axes
	\draw[->, thick] (-0.15, 0) -- (12, 0); % X-axis
	\draw[->, thick] (0, -0.15) -- (0, 7);  % Y-axis

	% 3. Draw Axis Ticks and Labels
	% X-axis ticks (at integer coordinates)
	\draw[thick] (2, -0.1) -- (2, 0.1) node[below=5pt] {$a_1^{(1)}$};
	\draw[thick] (6, -0.1) -- (6, 0.1) node[below=5pt] {$b_1^{(1)}$};

	\draw[thick] (4, -0.1) -- (4, 0.1) node[below=5pt] {$a_1^{(2)}$};
	\draw[thick] (7, -0.1) -- (7, 0.1) node[below=5pt] {$b_1^{(2)}$};

	\draw[thick] (9, -0.1) -- (9, 0.1) node[below=5pt] {$a_1^{(3)}$};
	\draw[thick] (10, -0.1) -- (10, 0.1) node[below=5pt] {$b_1^{(3)}$};

	% Y-axis ticks (at integer coordinates)
	\draw[thick] (-0.1, 3) -- (0.1, 3) node[left=5pt] {$a_2^{(1)}$};
	\draw[thick] (-0.1, 5) -- (0.1, 5) node[left=5pt] {$b_2^{(1)}$};

	\draw[thick] (-0.1, 2) -- (0.1, 2) node[left=5pt] {$a_2^{(2)}$};
	\draw[thick] (-0.1, 4) -- (0.1, 4) node[left=5pt] {$b_2^{(2)}$};

	% 4. Draw Rectangles (Edges exactly between grid points)
	% y-span is from 1.5 to 6.5
	\fill[Green, opacity=0.2] (2-0.5, 3-0.5) rectangle (6+0.5, 6-0.5);  % q1 bounding box
	\fill[gray, opacity=0.3] (4-0.5, 2-0.5) rectangle (7+0.5, 5-0.5);  % q1 bounding box
	\fill[gray, opacity=0.4] (10-0.5, 3-0.5) rectangle (12+0.5, 6-0.5);  % q1 bounding box

	% 5. Add Rectangle Area Labels
	\node at (2.5, 2) {$q_1$};
	\node at (5.5, 1) {$q_2$};
	\node at (11, 2) {$q_3$};
	
	% q1 marker
	\node[anchor=center, inner sep=2pt, Green] at (2,3) {$+C_1$};
	\node[anchor=center, inner sep=2pt, Green] at (7,3) {$-C_1$};
	\node[anchor=center, inner sep=2pt, Green] at (7,6) {$+C_1$};
	\node[anchor=center, inner sep=2pt, Green] at (2,6) {$-C_1$};

	% q2 marker
	\node[anchor=center, inner sep=2pt] at (4,2) {$+C_2$};
	\node[anchor=center, inner sep=2pt] at (8,2) {$-C_2$};
	\node[anchor=center, inner sep=2pt] at (8,5) {$+C_2$};
	\node[anchor=center, inner sep=2pt] at (4,5) {$-C_2$};

	% q3 updates
	\node[anchor=center, inner sep=2pt] at (10,3) {$+C_3$};
	%\node[anchor=center, inner sep=2pt] at (11,3) {$-C_3$};
	\node[anchor=center, inner sep=2pt] at (10,6) {$-C_3$};
	%\node[anchor=center, inner sep=2pt] at (11,6) {$+C_3$};
    \end{tikzpicture}
\end{frame}

\begin{frame}
    \vspace{.5cm}
    \begin{tikzpicture}
	\draw[help lines, color=gray!30, step=1cm] (0,0) grid (12, 7);

	% 2. Draw Axes
	\draw[->, thick] (-0.15, 0) -- (12, 0); % X-axis
	\draw[->, thick] (0, -0.15) -- (0, 7);  % Y-axis

	% 3. Draw Axis Ticks and Labels
	% X-axis ticks (at integer coordinates)
	\draw[thick] (2, -0.1) -- (2, 0.1) node[below=5pt] {$a_1^{(1)}$};
	\draw[thick] (6, -0.1) -- (6, 0.1) node[below=5pt] {$b_1^{(1)}$};

	\draw[thick] (4, -0.1) -- (4, 0.1) node[below=5pt] {$a_1^{(2)}$};
	\draw[thick] (7, -0.1) -- (7, 0.1) node[below=5pt] {$b_1^{(2)}$};

	\draw[thick] (9, -0.1) -- (9, 0.1) node[below=5pt] {$a_1^{(3)}$};
	\draw[thick] (10, -0.1) -- (10, 0.1) node[below=5pt] {$b_1^{(3)}$};

	% Y-axis ticks (at integer coordinates)
	\draw[thick] (-0.1, 3) -- (0.1, 3) node[left=5pt] {$a_2^{(1)}$};
	\draw[thick] (-0.1, 5) -- (0.1, 5) node[left=5pt] {$b_2^{(1)}$};

	\draw[thick] (-0.1, 2) -- (0.1, 2) node[left=5pt] {$a_2^{(2)}$};
	\draw[thick] (-0.1, 4) -- (0.1, 4) node[left=5pt] {$b_2^{(2)}$};

	% 4. Draw Rectangles (Edges exactly between grid points)
	% y-span is from 1.5 to 6.5
	\fill[gray, opacity=0.2] (2-0.5, 3-0.5) rectangle (6+0.5, 6-0.5);  % q1 bounding box
	\fill[gray, opacity=0.3] (4-0.5, 2-0.5) rectangle (7+0.5, 5-0.5);  % q1 bounding box
	\fill[gray, opacity=0.2, Green] (10-0.5, 3-0.5) rectangle (12+0.5, 6-0.5);  % q1 bounding box

	% 5. Add Rectangle Area Labels
	\node at (2.5, 2) {$q_1$};
	\node at (5.5, 1) {$q_2$};
	\node at (11, 2) {$q_3$};
	
	% q1 marker
	\node[anchor=center, inner sep=2pt] at (2,3) {$+C_1$};
	\node[anchor=center, inner sep=2pt] at (7,3) {$-C_1$};
	\node[anchor=center, inner sep=2pt] at (7,6) {$+C_1$};
	\node[anchor=center, inner sep=2pt] at (2,6) {$-C_1$};

	% q2 marker
	\node[anchor=center, inner sep=2pt] at (4,2) {$+C_2$};
	\node[anchor=center, inner sep=2pt] at (8,2) {$-C_2$};
	\node[anchor=center, inner sep=2pt] at (8,5) {$+C_2$};
	\node[anchor=center, inner sep=2pt] at (4,5) {$-C_2$};

	% q3 updates
	\node[anchor=center, inner sep=2pt, Green] at (10,3) {$+C_3$};
	%\node[anchor=center, inner sep=2pt] at (11,3) {$-C_3$};
	\node[anchor=center, inner sep=2pt, Green] at (10,6) {$-C_3$};
	%\node[anchor=center, inner sep=2pt] at (11,6) {$+C_3$};
    \end{tikzpicture}
\end{frame}

\begin{frame}
    \begin{definition*}{}
	Let \( D = \prod_{i=1}^d \llbracket n_i, m_i \rrbracket \subseteq \mathbb{Z}^d \) be a finite grid domain and \( q_1, \dots, q_L \) a series of constant update queries on \( D \) over the update ranges \( D_j = \prod_{i=1}^d \llbracket a_i^{(j)}, b_i^{(j)} \rrbracket \) with update values \( C_j \) for \( j \in \{1,\dots,L\} \).

	The \emph{difference array} \( \Delta[q_1,\dots,q_L] : D \to \mathbb{R} \) is defined as
	\begin{equation*}
	    \Delta[q_1,\dots,q_L] = \sum_{j=1}^L \sum_{p \in M_j} (-1)^{N_j(p)} \cdot C_j \cdot \delta_p,
	\end{equation*}
	where \( M_j = \prod_{i=1}^d \left( \left\{a_i^{(j)}\right\} \cup \left\{ b_i^{(j)} + 1 \mid b_i^{(j)} \neq m_i \right\} \right) \) and \( N_j(p) = \card\{ p_i \mid p_i = b_i^{(j)}+1 \} \) for \( j \in \{1,\dots,L\} \).
    \end{definition*}

    The notation \( \delta_p \) means the Dirac-delta given by 
    \begin{equation*}
	\delta_p(k) = \begin{cases} 1, & \text{if } k = p, \\ 0, & \text{else.} \end{cases}
    \end{equation*}
\end{frame}

\begin{frame}
    \begin{algorithmic}[1]
	\Statex\hspace{-2em} \textbf{Procedure:} \( \textsc{DiffArr}(D,q_1,\dots,q_Q) \)

	\Statex\hspace{-2em} \textbf{Input:} A finite grid domain \( D = \prod_{i=1}^d \llbracket n_i, m_i \rrbracket \subseteq \mathbb{Z}^d \) and a series of constant update queries \( q_1, \dots, q_Q \) on \( D \) over the update ranges \( D_j = \prod_{i=1}^d \llbracket a_i^{(j)}, b_i^{(j)} \rrbracket \) with update values \( C_j \) for \( 1 \leq j \leq Q \)

	\Statex\hspace{-2em} \textbf{Output:} The difference array \( \Delta[q_1,\dots,q_Q] \)

	\Statex

	\State \( d \gets \) Zero map \( d : D \to \mathbb{R} \) \hspace{2em}

	\For{ \( j \in \llbracket 1, Q \rrbracket \) }
	    \State \( M_j \gets \prod_{i=1}^d \left( \left\{a_i^{(j)}\right\} \cup \left\{ b_i^{(j)} + 1 \mid b_i^{(j)} \neq m_i \right\} \right) \)

	    \For{ \( p \in M_j \) }
		\State \( N(p) \gets \card\{ p_i \mid p_i = b_i^{(j)}+1 \} \)
		\State \( d(p) \gets d(p) + (-1)^{N(p)}\cdot C_j \)
	    \EndFor
	\EndFor

	\State \Return \( d \)
    \end{algorithmic}
\end{frame}

\begin{frame}
    \begin{proposition*}{}
	Let \( D = \prod_{i=1}^d \llbracket n_i, m_i \rrbracket \subseteq{Z}^d \) be a finite grid domain, \( N = \card(D) \) and let \( q_1, \dots, q_L \) be a series of constant update queries on \( D \) of arbitrary size. 
	The difference array algorithm computes the difference array \( \Delta[q_1,\dots,q_Q] \) in \( \mathcal{O}(N + L) \). 
    \end{proposition*}

    How does this formal notion of a difference array help us in computing the updated map \( F = f + \sum_{i=0}^L q_i \) more efficiently? 

    \vspace{.2cm}

    From the 1D/2D heuristic: Cumulatively summing over \( \Delta[q_1,\dots,q_L] \) should yield the total update \( \sum_{i=0}^L q_i \). 
\end{frame}

\begin{frame}{Roadmap: How does Discrete Calculus come into play?}
    \begin{itemize}
	\item We will see that a difference array \( \Delta[q_1,\dots,q_L] \) is the sum of discrete derivatives of the queries \( q_1, \dots, q_L \), that is 
	    \begin{equation*}
		\Delta[q_1,\dots,q_L] = \sum_{i=0}^L \Delta q_i. 
	    \end{equation*}
	\item Cumulatively summing over \( \Delta[q_1,\dots,q_L] \) will be the discrete analogon of integration and thus will be the inverse operation to discrete differentiation, that is 
	    \begin{equation*}
		\Sigma(\Delta[q_1,\dots,q_L]) = \sum_{i=0}^L \Sigma(\Delta q_i) = \sum_{i=0}^L q_i. 
	    \end{equation*}
    \end{itemize}
\end{frame}

\begin{frame}{Foundations of Discrete Calculus}
    \begin{definition*}{}
	The \emph{\( i \)-th backward difference} \( \Delta_i \) is the linear operator assigning to every discrete function \( f : D \subseteq \mathbb Z^d \to \mathbb R \) with \( D = \prod_{i=1}^d \llbracket n_i, m_i \rrbracket \subseteq \mathbb Z^d \) the function \( \Delta_i f: D \to \mathbb R \) given by 
	\begin{align*}
	    \Delta_i f(k_1, \dots, k_d) &:= f(k_1, \dots, k_{i-1}, k_i, k_{i+1}, \dots, k_d) \\
					& \quad\qquad - f(k_1, \dots, k_{i-1}, k_i-1, k_{i+1}, \dots, k_d)
	\end{align*}
	for all \( (k_1,\dots,k_d) \in D \) with \( k_i > n_i \) and 
	\begin{equation*}
	    \Delta_i f(k_1, \dots, k_d) := f(k_1, \dots, k_{i-1},k_i,k_{i+1}, \dots, k_d)
	\end{equation*}
	for all \( (k_1,\dots,k_d) \in D \) with \( k_i = n_i \). 
    \end{definition*}
\end{frame}

\begin{frame}
    \begin{definition*}{continued}
	For shorthand notation, we use 
	\begin{equation*}
	    \Delta := \Delta_1 \circ \cdots \circ \Delta_d
	\end{equation*}
	to denote the \emph{mixed backward difference} \( \Delta \). 
    \end{definition*}
\end{frame}

\begin{frame}
    \begin{definition*}{}
	The \emph{\( i \)-th (inclusive) prefix sum} \( \Sigma_i \) is the linear operator assigning to every discrete function \( f: D \subseteq \mathbb Z^d \to \mathbb R \) with \( D = \prod_{j=1}^d \llbracket n_i, m_i \rrbracket \) the function \( \Sigma_i f : D \to \mathbb R \) given by 
	\begin{equation*}
	    \Sigma_i f(k_1, \dots, k_d) := \sum_{t = n_i}^{k_i} f(k_1, \dots, k_{i-1}, t, k_{i+1}, \dots, k_d)
	\end{equation*}
	for all \( (k_1, \dots, k_d) \in D \).
	Again for shorthand notation, we  use
	\begin{equation*}
	    \Sigma := \Sigma_1 \circ \cdots \circ \Sigma_d
	\end{equation*}
	to denote the \emph{(inclusive) prefix sum} \( \Sigma \).
    \end{definition*}
\end{frame}

\begin{frame}
    It is not hard to show that backward differences and prefix sums commute, that is for every \( i \) and \( j \) we have \( \Sigma_i \circ \Delta_j = \Delta_j \circ \Sigma_i \). 

    \vspace{.4cm}

    In particular, \( \Sigma \circ \Delta = (\Sigma_1 \circ \Delta_1) \circ \cdots \circ (\Sigma_d \circ \Delta_d) =  \Delta \circ \Sigma \). 

    \begin{proposition*}{}
	The \( i \)-th backward difference \( \Delta_i \) and the \( i \)-th prefix sum \( \Sigma_i \) are inverse to each other for all \( i \). 
	That is, for every discrete function \( f : D \to \mathbb R \) with \( D = \prod_{i=1}^d \llbracket n_i, m_i \rrbracket \subseteq \mathbb Z^d \) we have
	\begin{equation*}
	    \Delta_i \circ \Sigma_i (f) = f \quad \text{and} \quad \Sigma_i \circ \Delta_i (f) = f
	\end{equation*}
	for all \( i \in \{1,\dots,d\} \). 
	In particular, the mixed backward difference \( \Delta = \Delta_1 \circ \cdots \circ \Delta_d \) and the prefix sum \( \Sigma = \Sigma_1 \circ \cdots \circ \Sigma_d \) are inverse to each other. 
    \end{proposition*}

    Another important property of prefix sums is that they can be computed efficiently.
\end{frame}

\begin{frame}
    \begin{algorithmic}[1]
	\Statex\hspace{-2em} \textbf{Procedure:} \textsc{PrefixSum}
	\Statex\hspace{-2em} \textbf{Input:} A discrete map \( f : D \to \mathbb R \) with \( D = \prod_{i=1}^d \llbracket n_i, m_i \rrbracket \subseteq \mathbb{Z}^d \)

	\Statex\hspace{-2em} \textbf{Output:} The prefix sum \( \Sigma(f) : D \to \mathbb{R} \)

	\Statex 

	\State \( S \gets \) The discrete map \( f : D \to \mathbb{R} \) \hspace{2em}

	\For{ \( j \in \llbracket 1, d \rrbracket \) }
    \For{ \( (k_1,\dots,k_{j-1}, k_{j+1}, \dots, k_d) \in \prod_{\substack{i=1\\i\neq j}}^d \llbracket n_i, m_i \rrbracket \) }
		\For{ \( t \in \llbracket n_j+1, m_j \rrbracket \) }
			\State \( S(k_1,\dots,k_{j-1}, t, k_{j+1}, \dots,k_d) \gets S(k_1, \dots, k_{j-1}, t-1, k_{j+1}, \dots, k_d) \) \\ \( \phantom{S(k_1,\dots,k_{j-1}, t, k_{j+1}, \dots,k_d) \gets} + S(k_1,\dots,k_{j-1},t,k_{j+1},\dots,k_d) \)
		\EndFor
	    \EndFor
	\EndFor

	\State \Return \( S \)
    \end{algorithmic}
\end{frame}

\begin{frame}
    \begin{proposition*}{}
	Let \( f : D \to \mathbb R \) be a discrete map with \( D = \prod_{i=1}^d \llbracket n_i, m_i \rrbracket \subseteq \mathbb{Z}^d \).
	The \mbox{\textsc{PrefixSum}} algorithm computes the prefix sum \( \Sigma f : D \to \mathbb{R} \) in \( \mathcal{O}(N) \) with \( N = \card(D) \). 
    \end{proposition*}
\end{frame}

\begin{frame}{Indicator functions and Heaviside step functions}
    \begin{definition*}{}{}
	Let \( p = (p_1, \dots, p_d) \in \mathbb Z^d \). 
	The \emph{Heaviside step function} \( H_p : \mathbb Z^d \to \mathbb R \) is the indicator function given by 
	\begin{equation*}
	    H_p(k_1, \dots, k_d) := \begin{cases}
		1, & \text{if } \forall 1 \leq i \leq d : k_i \geq p_i, \\
		0, & \text{else}
	    \end{cases}
	\end{equation*}
	for \( (k_1,\dots,k_d) \in \mathbb{Z}^d \). 
    \end{definition*}

    Heaviside step functions are specific indicator functions with very simple mixed backward differences.
\end{frame}

\begin{frame}
    \begin{proposition*}{}
	Let \( D = \prod_{i=1}^d \llbracket n_i, m_i \rrbracket \subseteq \mathbb Z^d \) be a finite grid domain and \( p = (p_1, \dots, p_d) \in D \). 
	The mixed backward difference \( \Delta H_p \) of the Heaviside step function \( H_p : D \to \mathbb R \) is given by 
	\begin{equation*}
	    \Delta H_p(k_1, \dots, k_d) = \delta_p := \begin{cases}
		1, & \text{if } \forall 1 \leq i \leq d : k_i = p_i, \\
		0, & \text{else}.
	    \end{cases}
	\end{equation*}
    \end{proposition*}
    \begin{proof}
	\begin{itemize}
	    \item It is straight forward to see that \( H_p = \prod_{i=1}^d H_{p_i} \)
	    \item One can show that for such products we have \( \Delta H_p = \prod_{i=1}^d \Delta H_{p_i} \)
	    \item Show \( \Delta H_{p_i} = \delta_{p_i} \) by a simple case distinction
	    \item Finally, \( \Delta H_p = \prod_{i=1}^d \delta_{p_i} = \delta_{(p_1,\dots,p_d)} \) \qedhere
	\end{itemize}
    \end{proof}
\end{frame}

\begin{frame}
    \begin{proposition*}{}
	Let \( D = \prod_{i=1}^d \llbracket a_i, b_i \rrbracket \subseteq \mathbb Z^d \) be a finite grid domain. 
	Then, the indicator function \( \mathbb{1}_D : \mathbb Z^d \to \mathbb R \) can be written as 
	\begin{equation*}
	    \mathbb{1}_D = \sum_{p \in \prod_{i=1}^d \{a_i,b_i+1\}} (-1)^{N(p)} H_p,
	\end{equation*}
	where \( N(p) = \card\{ p_i \mid p_i = b_i+1 \} \). 
    \end{proposition*}
    \begin{proof}
	\begin{itemize}
	    \item It is not hard to see that \( \mathbb{1}_D(k_1,\dots,k_d) = \prod_{i=1}^d \mathbb{1}_{\llbracket a_i, b_i \rrbracket}(k_i) \)
	\item By a simple case distiction we see that \( \mathbb{1}_{\llbracket a_i, b_i \rrbracket}(k_i) = H_{a_i}(k_i) - H_{b_i+1}(k_i) \) 
	    \item Hence, \( \mathbb{1}_D(k_1,\dots,k_d) = \prod_{i=1}^d \left( H_{a_i}(k_i) - H_{b_i+1}(k_i) \right) \)
	    \item Using the general distribution formula for products yields 
		    \begin{align*}
			\prod_{i=1}^d \left( H_{a_i}(k_i) - H_{b_i+1}(k_i) \right) = \sum_{J \subseteq I} (-1)^{\card(J)} \Bigg( \prod_{i \in I \setminus J} H_{a_i}(k_i) \Bigg) \Bigg( \prod_{j \in J} H_{b_j+1}(k_j) \Bigg),
		    \end{align*}
		    where \( I = \{1,\dots,d\} \). \qedhere
	\end{itemize}
    \end{proof}
\end{frame}

\begin{frame}
    \begin{corollary*}{}
	Let \( D = \prod_{i=1}^d \llbracket n_i, m_i \rrbracket \subseteq \mathbb{Z}^d \) be a finite grid domain and let \( q : D \to \mathbb R \) be a constant update query on \( D \) with update range \( D_q = \prod_{i=1}^d \llbracket a_i, b_i \rrbracket \) and update value \( C \in \mathbb{R} \). 
	Then, 
	\begin{equation*}
	    q = \sum_{p \in M} (-1)^{N(p)} \cdot C \cdot H_p,
	\end{equation*}
	where \( M = \prod_{i=1}^d \left( \left\{a_i\right\} \cup \left\{ b_i + 1 \mid b_i \neq m_i \right\} \right) \) and \( N(p) = \card\{ p_i \mid p_i = b_i+1 \} \). 
	In particular, the mixed backward difference \( \Delta q \) of \( q \) can be expressed as 
	\begin{equation*}
	    \Delta q = \sum_{p \in M} (-1)^{N(p)} \cdot C \cdot \delta_p. 
	\end{equation*}
    \end{corollary*}
    \begin{proof}
	\begin{itemize}
	    \item By definition \( q = C \cdot \mathbb{1}_{D_q} \) 

	    \item Writing the indicator function by Heaviside step functions yields the first claim

	    \item The second claim follows by the linearity of \( \Delta \) and as previously shown \( \Delta H_p = \delta_p \)  \qedhere
	\end{itemize}
    \end{proof}
\end{frame}

\begin{frame}
    Recall the formal definition of difference arrays: 

    \begin{definition*}{}
	    Let \( D = \prod_{i=1}^d \llbracket n_i, m_i \rrbracket \subseteq \mathbb{Z}^d \) be a finite grid domain and \( q_1, \dots, q_L \) a series of constant update queries on \( D \). 
	    The \emph{difference array} \( \Delta[q_1,\dots,q_Q] : D \to \mathbb{R} \) is defined as
	\begin{equation*}
	    \Delta[q_1,\dots,q_L] = \sum_{j=1}^L \sum_{p \in M_j} (-1)^{N_j(p)} \cdot C_j \cdot \delta_p,
	\end{equation*}
	where \( M_j = \prod_{i=1}^d \left( \left\{a_i^{(j)}\right\} \cup \left\{ b_i^{(j)} + 1 \mid b_i^{(j)} \neq m_i \right\} \right) \) and \( N_j(p) = \card\{ p_i \mid p_i = b_i^{(j)}+1 \} \) for \( 1 \leq j \leq Q \).
    \end{definition*}

    \begin{itemize}
	\item With the corollary on the last slide we see that 
	    \begin{equation*}
		\Delta q_j = \sum_{p \in M_j} (-1)^{N_j(p)} \cdot C_j \cdot \delta_p
	    \end{equation*}
	    for all \( j \in \{1,\dots,L\} \). 
    \end{itemize}
\end{frame}

\begin{frame}
    \begin{itemize}
	\item Hence, we indeed have 
	\begin{equation*}
	    \Delta[q_1,\dots,q_L] = \sum_{j=1}^L \Delta q_j. 
	\end{equation*}
	\item The prefix sum and the mixed backwards difference are inverse to each other, thus 
	\begin{equation*}
	    \Sigma(\Delta[q_1,\dots,q_L]) = \sum_{j=1}^L \Sigma(\Delta q_j) = \sum_{j=1}^L q_j.
	\end{equation*}
	
	\item We have seen that \( \Delta[q_1,\dots,q_L] \) can be computed in \( \mathcal{O}(N + L) \). 

	\item We have seen that the prefix sum \( \Sigma(\Delta[q_1,\dots,q_L]) \) can be computed in \( \mathcal{O}(N) \). 

	\item The addition \( f + \Sigma(\Delta[q_1,\dots,q_l]) \) has time complexity \( \mathcal{O}(N) \). 

	\item In total: Computing \( f + \Sigma(\Delta[q_1,\dots,q_l]) \) can be done in \( \mathcal{O}(N + L) \). 
    \end{itemize}
\end{frame}

\begin{frame}
    \begin{algorithmic}[1]
	\Statex\hspace{-2em} \textbf{Input:} A discrete map \( f : D \to \mathbb R \) with \( D = \prod_{i=1}^d \llbracket n_i, m_i \rrbracket \) and a series of constant update queries \( q_1, \dots, q_Q \) on \( D \) 

	\Statex\hspace{-2em} \textbf{Output:} The updated map \( F = f + \sum_{i=1}^Q q_i \)

	\Statex

	\State \( d \gets \Call{DiffArr}{D,q_1,\dots,q_Q} \)

	\State \(\Sigma d \gets \Call{PrefixSum}{d} \)

	\State \( F \gets f + \Sigma d \) 

	\State \Return \( F \)
    \end{algorithmic}

    \begin{proposition*}{}
	Let \( f : D \subseteq \mathbb{Z}^d \to \mathbb{R} \) be a discrete function with \( D = \prod_{i=1}^d \llbracket n_i, m_i \rrbracket \) and \( N = \card(D) \), and let \( q_1, \dots, q_L \) be a series of constant update queries on \( D \) of arbitrary size. 
	Assuming that for each \( k \in D \) the function value \( f(k) \) can be obtained in constant time, the algorithm above computes the updated map \( F = f + \sum_{i=1}^L q_i \) in \( \mathcal{O}(N + L) \). 
    \end{proposition*}
\end{frame}

\begin{frame}{Flat update queries}
    \vfill

    \begin{tikzpicture}
	\draw[help lines, color=gray!30, step=1cm] (0,0) grid (12, 4);

	% 2. Draw Axes
	\draw[->, thick] (-0.15, 0) -- (12, 0); % X-axis
	\draw[->, thick] (0, -0.15) -- (0, 4);  % Y-axis

	% 3. Draw Axis Ticks and Labels
	% X-axis ticks (at integer coordinates)
	\draw[thick] (2, -0.1) -- (2, 0.1) node[below=5pt] {$a_1^{(1)}$};

	\draw[thick] (4, -0.1) -- (4, 0.1) node[below=5pt] {$a_1^{(2)}$};

	\draw[thick] (6, -0.1) -- (6, 0.1) node[below=5pt] {$a_1^{(3)}$};
	\draw[thick] (10, -0.1) -- (10, 0.1) node[below=5pt] {$b_1^{(3)}$};

	% Y-axis ticks (at integer coordinates)
	\draw[thick] (-0.1, 1) -- (0.1, 1) node[left=5pt] {$a_2^{(1)}$};
	\draw[thick] (-0.1, 3) -- (0.1, 3) node[left=5pt] {$b_2^{(1)}$};

	% 4. Draw Rectangles (Edges exactly between grid points)
	% y-span is from 1.5 to 6.5
	\draw[thick, gray!20, fill=gray, fill opacity=0.2] (2-0.5, 1-0.5) rectangle (2+0.5, 3+0.5);  % q1 bounding box
	\draw[thick, gray!30, fill=gray, fill opacity=0.3] (4-0.5, 1-0.5) rectangle (4+0.5, 3+0.5);  % q1 bounding box
	\draw[thick, gray!40, fill=gray, fill opacity=0.4] (6-0.5, 1-0.5) rectangle (10+0.5, 1+0.5);  % q1 bounding box

	% 5. Add Rectangle Area Labels
	\node at (1.25, 2) {$q_1$};
	\node at (3.25, 2) {$q_2$};
	\node at (5.25, 1) {$q_3$};
	
	% q1 updates
	\node[anchor=center, inner sep=2pt] at (2,1) {$+C_1$};
	\node[anchor=center, inner sep=2pt] at (2,2) {$+C_1$};
	\node[anchor=center, inner sep=2pt] at (2,3) {$+C_1$};

	% q2 updates
	\node[anchor=center, inner sep=2pt] at (4,1) {$+C_2$};
	\node[anchor=center, inner sep=2pt] at (4,2) {$+C_2$};
	\node[anchor=center, inner sep=2pt] at (4,3) {$+C_2$};

	% q3 updates
	\node[anchor=center, inner sep=2pt] at (6,1) {$+C_3$};
	\node[anchor=center, inner sep=2pt] at (7,1) {$+C_3$};
	\node[anchor=center, inner sep=2pt] at (8,1) {$+C_3$};
	\node[anchor=center, inner sep=2pt] at (9,1) {$+C_3$};
	\node[anchor=center, inner sep=2pt] at (10,1) {$+C_3$};
    \end{tikzpicture}

    \vfil
\end{frame}

\begin{frame}
    \vfill

    \begin{tikzpicture}
	\draw[help lines, color=gray!30, step=1cm] (0,0) grid (12, 5);

	% 2. Draw Axes
	\draw[->, thick] (-0.15, 0) -- (12, 0); % X-axis
	\draw[->, thick] (0, -0.15) -- (0, 5);  % Y-axis

	% 3. Draw Axis Ticks and Labels
	% X-axis ticks (at integer coordinates)
	\draw[thick] (2, -0.1) -- (2, 0.1) node[below=5pt] {$a_1^{(1)}$};

	\draw[thick] (4, -0.1) -- (4, 0.1) node[below=5pt] {$a_1^{(2)}$};

	\draw[thick] (6, -0.1) -- (6, 0.1) node[below=5pt] {$a_1^{(3)}$};
	\draw[thick] (10, -0.1) -- (10, 0.1) node[below=5pt] {$b_1^{(3)}$};

	% Y-axis ticks (at integer coordinates)
	\draw[thick] (-0.1, 1) -- (0.1, 1) node[left=5pt] {$a_2^{(1)}$};
	\draw[thick] (-0.1, 3) -- (0.1, 3) node[left=5pt] {$b_2^{(1)}$};

	% 4. Draw Rectangles (Edges exactly between grid points)
	% y-span is from 1.5 to 6.5
	\draw[thick, gray!20, fill=gray, fill opacity=0.2] (2-0.5, 1-0.5) rectangle (2+0.5, 3+0.5);  % q1 bounding box
	\draw[thick, gray!30, fill=gray, fill opacity=0.3] (4-0.5, 1-0.5) rectangle (4+0.5, 3+0.5);  % q1 bounding box
	\draw[thick, gray!40, fill=gray, fill opacity=0.4] (6-0.5, 1-0.5) rectangle (10+0.5, 1+0.5);  % q1 bounding box

	% 5. Add Rectangle Area Labels
	\node at (1.25, 2) {$q_1$};
	\node at (3.25, 2) {$q_2$};
	\node at (8, 1.75) {$q_3$};
	
	% q1 markers
	\node[anchor=center, inner sep=2pt] at (2,1) {$+C_1$};
	\node[anchor=center, inner sep=2pt] at (3,1) {$-C_1$};
	\node[anchor=center, inner sep=2pt] at (2,4) {$-C_1$};
	\node[anchor=center, inner sep=2pt] at (3,4) {$+C_1$};

	% q2 updates
	\node[anchor=center, inner sep=2pt] at (4,1) {$+C_2$};
	\node[anchor=center, inner sep=2pt] at (5,1) {$-C_2$};
	\node[anchor=center, inner sep=2pt] at (4,4) {$-C_2$};
	\node[anchor=center, inner sep=2pt] at (5,4) {$+C_2$};

	% q3 updates
	\node[anchor=center, inner sep=2pt] at (6,1) {$+C_3$};
	\node[anchor=center, inner sep=2pt] at (6,2) {$-C_3$};
	\node[anchor=center, inner sep=2pt] at (11,1) {$-C_3$};
	\node[anchor=center, inner sep=2pt] at (11,2) {$+C_3$};
    \end{tikzpicture}

    \vfil
\end{frame}

\begin{frame}
    \vfill

    \begin{tikzpicture}
	\draw[help lines, color=gray!30, step=1cm] (0,0) grid (12, 5);

	% 2. Draw Axes
	\draw[->, thick] (-0.15, 0) -- (12, 0); % X-axis
	\draw[->, thick] (0, -0.15) -- (0, 5);  % Y-axis

	% 3. Draw Axis Ticks and Labels
	% X-axis ticks (at integer coordinates)
	\draw[thick] (2, -0.1) -- (2, 0.1) node[below=5pt] {$a_1^{(1)}$};

	\draw[thick] (4, -0.1) -- (4, 0.1) node[below=5pt] {$a_1^{(2)}$};

	\draw[thick] (6, -0.1) -- (6, 0.1) node[below=5pt] {$a_1^{(3)}$};
	\draw[thick] (10, -0.1) -- (10, 0.1) node[below=5pt] {$b_1^{(3)}$};

	% Y-axis ticks (at integer coordinates)
	\draw[thick] (-0.1, 1) -- (0.1, 1) node[left=5pt] {$a_2^{(1)}$};
	\draw[thick] (-0.1, 3) -- (0.1, 3) node[left=5pt] {$b_2^{(1)}$};

	% 4. Draw Rectangles (Edges exactly between grid points)
	% y-span is from 1.5 to 6.5
	\draw[thick, gray!20, fill=gray, fill opacity=0.2] (2-0.5, 1-0.5) rectangle (2+0.5, 3+0.5);  % q1 bounding box
	\draw[thick, gray!30, fill=gray, fill opacity=0.3] (4-0.5, 1-0.5) rectangle (4+0.5, 3+0.5);  % q1 bounding box
	\draw[thick, gray!40, fill=gray, fill opacity=0.4] (6-0.5, 1-0.5) rectangle (10+0.5, 1+0.5);  % q1 bounding box

	% 5. Add Rectangle Area Labels
	\node at (1.25, 2) {$q_1$};
	\node at (3.25, 2) {$q_2$};
	\node at (8, 1.75) {$q_3$};
	
	% q1 markers
	\node[anchor=center, inner sep=2pt] at (2,1) {$+C_1$};
	\node[anchor=center, inner sep=2pt] at (2,4) {$-C_1$};

	% q2 updates
	\node[anchor=center, inner sep=2pt] at (4,1) {$+C_2$};
	\node[anchor=center, inner sep=2pt] at (4,4) {$+C_2$};

	% q3 updates
	\node[anchor=center, inner sep=2pt] at (6,1) {$+C_3$};
	\node[anchor=center, inner sep=2pt] at (11,1) {$-C_3$};
    \end{tikzpicture}

    \vfil
\end{frame}

\begin{frame}{Consecutive update queries}
    \vfill

    \begin{tikzpicture}
	% Define custom colors to match the screenshot
	\definecolor{myyellow}{RGB}{230, 180, 0}
	\definecolor{mygreen}{RGB}{30, 160, 30}

	% Coords
	\def\ax{2}
	\def\ay{2}
	\def\widthone{3}
	\def\widthtwo{2}
	\def\widththree{4}
	\def\height{4}

	\pgfmathtruncatemacro{\xstartone}{\ax}
	\pgfmathtruncatemacro{\xendone}{\ax + \widthone - 1}
	\pgfmathtruncatemacro{\xstarttwo}{\ax + \widthone}
	\pgfmathtruncatemacro{\xendtwo}{\ax + \widthone + \widthtwo - 1}
	\pgfmathtruncatemacro{\xstartthree}{\ax + \widthone + \widthtwo}
	\pgfmathtruncatemacro{\xendthree}{\ax + \widthone + \widthtwo + \widththree - 1}

	\pgfmathtruncatemacro{\ystart}{\ay}
	\pgfmathtruncatemacro{\yend}{\ay+\height-1}

	% 1. Draw Background Grid (Gray Dots)
	% Limits chosen to cover the entire drawing area with a margin
	%\foreach \x in {0, ..., 13} {
	%    \foreach \y in {0, ..., 7} {
	%        \fill (\x,\y) circle (0.8pt);
	%    }
	%}

	\draw[help lines, color=gray!30, step=1cm] (0,0) grid (\xendthree+2, \yend+1);

	% 2. Draw Axes
	\draw[->, thick] (-0.15, 0) -- (\xendthree+2, 0); % X-axis
	\draw[->, thick] (0, -0.15) -- (0, \yend+1);  % Y-axis

	% 3. Draw Axis Ticks and Labels
	% X-axis ticks (at integer coordinates)
	\draw[thick] (\xstartone, -0.1) -- (\xstartone, 0.1) node[below=5pt] {$a_1^{(1)}$};
	\draw[thick] (\xendone, -0.1) -- (\xendone, 0.1) node[below=5pt] {$b_1^{(1)}$};
	\draw[thick] (\xstarttwo, -0.1) -- (\xstarttwo, 0.1) node[below=5pt] {$a_1^{(2)}$};
	\draw[thick] (\xendtwo, -0.1) -- (\xendtwo, 0.1) node[below=5pt] {$b_1^{(2)}$};
	\draw[thick] (\xstartthree, -0.1) -- (\xstartthree, 0.1) node[below=5pt] {$a_1^{(3)}$};
	\draw[thick] (\xendthree, -0.1) -- (\xendthree, 0.1) node[below=5pt] {$b_1^{(3)}$};

	% Y-axis ticks (at integer coordinates)
	\draw[thick] (-0.1, \ystart) -- (0.1, \ystart) node[left=5pt] {$a_2^{(1)}$};
	\draw[thick] (-0.1, \yend) -- (0.1, \yend) node[left=5pt] {$b_2^{(1)}$};

	% 4. Draw Rectangles (Edges exactly between grid points)
	% y-span is from 1.5 to 6.5
	\draw[thick, gray!20, fill=gray, fill opacity=0.2] (\xstartone-0.5, \ystart-0.5) rectangle (\xendone+0.5, \yend+0.5);  % q1 bounding box
	\draw[thick, gray!30, fill=gray, fill opacity=0.3] (\xstarttwo-0.5, \ystart-0.5) rectangle (\xendtwo+0.5, \yend+0.5);  % q2 bounding box
	\draw[thick, gray!40, fill=gray, fill opacity=0.4] (\xstartthree-0.5, \ystart-0.5) rectangle (\xendthree+0.5, \yend+0.5); % q3 bounding box

	% 5. Add Rectangle Area Labels
	\node at (\xstartone + \widthone*0.5 - 0.5, 1) {$q_1$};
	\node at (\xstarttwo + \widthtwo*0.5 -0.5, 1) {$q_2$};
	\node at (\xstartthree + \widththree*0.5 -0.5, 1) {$q_3$};

	% q1 updates
	\foreach \x in {\xstartone, ..., \xendone} {
	    \foreach \y in {\ystart, ..., \yend} {
		\node[anchor=center, inner sep=2pt] at (\x,\y) {$+C_1$};
	    }
	}

	% q2 updates
	\foreach \x in {\xstarttwo, ..., \xendtwo} {
	    \foreach \y in {\ystart, ..., \yend} {
		\node[anchor=center, inner sep=2pt] at (\x,\y) {$+C_2$};
	    }
	}

	% q3 updates
	\foreach \x in {\xstartthree, ..., \xendthree} {
	    \foreach \y in {\ystart, ..., \yend} {
		\node[anchor=center, inner sep=2pt] at (\x,\y) {$+C_3$};
	    }
	}
    \end{tikzpicture}

    \vfil
\end{frame}

\begin{frame}
    \vfill

    \begin{tikzpicture}
	% Define custom colors to match the screenshot
	\definecolor{myyellow}{RGB}{230, 180, 0}
	\definecolor{mygreen}{RGB}{30, 160, 30}

	% Coords
	\def\ax{2}
	\def\ay{2}
	\def\widthone{3}
	\def\widthtwo{2}
	\def\widththree{4}
	\def\height{4}

	\pgfmathtruncatemacro{\xstartone}{\ax}
	\pgfmathtruncatemacro{\xendone}{\ax + \widthone - 1}
	\pgfmathtruncatemacro{\xstarttwo}{\ax + \widthone}
	\pgfmathtruncatemacro{\xendtwo}{\ax + \widthone + \widthtwo - 1}
	\pgfmathtruncatemacro{\xstartthree}{\ax + \widthone + \widthtwo}
	\pgfmathtruncatemacro{\xendthree}{\ax + \widthone + \widthtwo + \widththree - 1}

	\pgfmathtruncatemacro{\ystart}{\ay}
	\pgfmathtruncatemacro{\yend}{\ay+\height-1}

	% 1. Draw Background Grid (Gray Dots)
	% Limits chosen to cover the entire drawing area with a margin
	%\foreach \x in {0, ..., 13} {
	%    \foreach \y in {0, ..., 7} {
	%        \fill (\x,\y) circle (0.8pt);
	%    }
	%}

	\draw[help lines, color=gray!30, step=1cm] (0,0) grid (\xendthree+2, \yend+2);

	% 2. Draw Axes
	\draw[->, thick] (-0.15, 0) -- (\xendthree+2, 0); % X-axis
	\draw[->, thick] (0, -0.15) -- (0, \yend+2);  % Y-axis

	% 3. Draw Axis Ticks and Labels
	% X-axis ticks (at integer coordinates)
	\draw[thick] (\xstartone, -0.1) -- (\xstartone, 0.1) node[below=5pt] {$a_1^{(1)}$};
	\draw[thick] (\xendone, -0.1) -- (\xendone, 0.1) node[below=5pt] {$b_1^{(1)}$};
	\draw[thick] (\xstarttwo, -0.1) -- (\xstarttwo, 0.1) node[below=5pt] {$a_1^{(2)}$};
	\draw[thick] (\xendtwo, -0.1) -- (\xendtwo, 0.1) node[below=5pt] {$b_1^{(2)}$};
	\draw[thick] (\xstartthree, -0.1) -- (\xstartthree, 0.1) node[below=5pt] {$a_1^{(3)}$};
	\draw[thick] (\xendthree, -0.1) -- (\xendthree, 0.1) node[below=5pt] {$b_1^{(3)}$};

	% Y-axis ticks (at integer coordinates)
	\draw[thick] (-0.1, \ystart) -- (0.1, \ystart) node[left=5pt] {$a_2^{(1)}$};
	\draw[thick] (-0.1, \yend) -- (0.1, \yend) node[left=5pt] {$b_2^{(1)}$};

	% 4. Draw Rectangles (Edges exactly between grid points)
	% y-span is from 1.5 to 6.5
	\draw[thick, gray!20, fill=gray, fill opacity=0.2] (\xstartone-0.5, \ystart-0.5) rectangle (\xendone+0.5, \yend+0.5);  % q1 bounding box
	\draw[thick, gray!30, fill=gray, fill opacity=0.3] (\xstarttwo-0.5, \ystart-0.5) rectangle (\xendtwo+0.5, \yend+0.5);  % q2 bounding box
	\draw[thick, gray!40, fill=gray, fill opacity=0.4] (\xstartthree-0.5, \ystart-0.5) rectangle (\xendthree+0.5, \yend+0.5); % q3 bounding box

	% 5. Add Rectangle Area Labels
	\node at (\xstartone + \widthone*0.5 - 0.5, 1) {$q_1$};
	\node at (\xstarttwo + \widthtwo*0.5 -0.5, 1) {$q_2$};
	\node at (\xstartthree + \widththree*0.5 -0.5, 1) {$q_3$};

	% 6. Add Corner/Boundary Labels
	% q1 Left corners (Black)
	\node[anchor=center, inner sep=2pt] at (\xstartone, \yend+1) {$-C_1$};
	\node[anchor=center, inner sep=2pt] at (\xstartone, \ystart) {$+C_1$};

	% q1/q2 Boundary labels
	\node[anchor=center, align=center, inner sep=2pt] 
	    at (\xstarttwo, \yend+1) {$+C_1$\\[-0.5ex]$-C_2$};
	\node[anchor=center, align=left, inner sep=2pt] 
	    at (\xstarttwo, \ystart) {$-C_1$\\[-0.5ex]$+C_2$};

	% q2/q3 Boundary labels
	\node[anchor=center, align=center, inner sep=2pt] 
	    at (\xstartthree, \yend+1) {$+C_2$\\[-0.5ex]$-C_3$};
	\node[anchor=center, align=left, inner sep=2pt] 
	    at (\xstartthree, \ystart) {$-C_2$\\[-0.5ex]$+C_3$};

	% q3 Right corners (Black)
	\node[anchor=center, inner sep=2pt] at (\xendthree+1, \yend+1) {$+C_3$};
	\node[anchor=center, inner sep=4pt] at (\xendthree+1, \ystart) {$-C_3$};
    \end{tikzpicture}

    \vfil
\end{frame}

\end{document}
